\chapter*{Conclusion}
\addcontentsline{toc}{chapter}{Conclusion}

%compared to our supposition taht memory is main bottleneck,
%runtime is a major issue
%
%in case of sg, need to increse the ram or decrase the number of workers
%
%in case of pinn, need to run longer, and larger arch need even longer
%
%
%PINNS
%
%time hungry
%non aparent what arch to use
%
%SG
%spatial adaptivity while retaining combi tech - still a an active fielf of reerach.
%\cite{discoadapt}
%
%
%DOG similar?
%when who better?
%low dim sg faster?
\begin{figure}
    \centering
    \includegraphics[width=\linewidth]{figures/threshold_cost_comparison_rel_l2.png}
    \caption{PINN vs combination technique performance in higher dimensions.
    In the context of PINNs, the degrees of freedom denote the number of trainable parameters.
    For the combination technique, it counts the total number of grid points in all component grids.}
    \label{fig:comparison_epilog}
\end{figure}


Both the combination technique and physics-informed neural networks are clearly influenced by the curse of dimensionality.
However, the effect is quite different for each method.
For sparse grids, the bottleneck is mainly combinatorial growth in the number of component grids and the corresponding memory cost.
For PINNs, the bottleneck is less abrupt, but it appears through the steadily growing need for larger networks, and longer training horizons.

The sparse-grid combination technique provides a strong method in low to moderate dimensions,
but its computational and memory cost rises quickly with the dimension and the level required to maintain an acceptable error tolerance.
Crucially, as the level or dimension increases, the memory needs of each worker rise as well.
Thus, in higher dimensions, somewhat contrary to the parallel advantage offered by the combination technique,
one must either dramatically increase the available memory or reduce the number of workers to fit the
growing number of component grids and grid points in memory.

Lastly, in \ref{fig:comparison_epilog} we compare the degrees of freedom needed by the combination technique and physics-informed neural networks to achieve a given accuracy.
In lower dimensions, sparse grids offer much larger transparency, accuracy, and speed.
As the dimension grows, however, the PINN approach appears to degrade more slowly,
even though this comes at the cost of longer training windows and larger models.
%One important difference, however, is that the sparse-grid solver can be treated as a relatively direct numerical method once its discretization is fixed.
%By contrast, the PINN approach requires substantially more tuning, especially in the first few iterations.

In future work, an interesting avenue for further study would be to explore is a spatially adaptive version of the combiantion technique.
Though harder on implementation, it promises to offer an even better compromise between approximation accuracy and degrees of freedom.